\chapter{资源与作用域}
\label{ch:08}

前面两章讲了「值」和「引用」，但一直没回答一个很现实的问题：

\begin{tip}{提示}
没有垃圾回收，那内存、文件句柄、网络连接这些东西，什么时候释放？谁负责？
\end{tip}

ChaLi 的答案很短：离开作用域的时候，自动释放。

这一章讲这套机制怎么用、为什么它不出错，以及堆对象该怎么管。

\section{作用域就是一对花括号}

先说清「作用域」在这门语言里指什么。

每一对 \texttt{\{ \}} 划出一个作用域。 在那个花括号里声明的东西，出了花括号就不存在了：

\par
\medskip
\begin{lstlisting}[style=chali]
void main()
{
    int64 outer = 1;
    {
        int64 inner = 2;
        println("里面能看到 {} 和 {}", outer, inner);
    }
    println("外面只剩 {}", outer);

}
\end{lstlisting}
\par
\medskip

\par
\medskip
\begin{lstlisting}[style=out]
里面能看到 1 和 2
外面只剩 1
\end{lstlisting}
\par
\medskip

这里面 \texttt{inner} 只活在里层那对大括号里。出了它，\texttt{inner} 就没了——想在外面用 \texttt{inner} 是编译错误，不是「值是 0」。

作用域有好几种来源：

{\small\setlength{\tabcolsep}{4pt}
\begin{longtable}{>{\raggedright\arraybackslash}p{\dimexpr 0.296\linewidth-2\tabcolsep\relax}>{\raggedright\arraybackslash}p{\dimexpr 0.684\linewidth-2\tabcolsep\relax}}
\toprule
来源 & 例子 \\
\midrule
\endhead
一对花括号 & \texttt{\{ int64 x =\allowbreak{} 1; \}} \\
函数体 & \texttt{int64 add(\allowbreak{}.\allowbreak{}.\allowbreak{}.\allowbreak{})\allowbreak{} \{ .\allowbreak{}.\allowbreak{}.\allowbreak{} \}} —— 整个函数是一个作用域 \\
\texttt{if} / \texttt{for} / \texttt{while} 的主体 & \texttt{for (\allowbreak{}int64 k in std:\allowbreak{}:\allowbreak{}lang:\allowbreak{}:\allowbreak{}range(\allowbreak{}0,\allowbreak{} 3)\allowbreak{})\allowbreak{} \{ \}} 里的 \texttt{k} 只活在循环里 \\
块 & \texttt{main} 本身也是块 \\
\bottomrule
\end{longtable}
}

注意最后一条：\texttt{main} 也是个作用域。 所以 \texttt{main} 里声明的东西，程序结束时会被清理。

\section{\texttt{Drop}：离开作用域时自动调用}

给一个结构体实现 \texttt{Drop}，它的 \texttt{drop(\allowbreak{})\allowbreak{}} 就会在这个值离开作用域时被调用：

\par
\medskip
\begin{lstlisting}[style=chali]
struct Resource: Drop
{
    String name;
    void drop()
    {
        println("drop {}", self.name);
    }
};

void main()
{
    {
        Resource r;
        r.name = "conn";
        println("using");
    }
    println("after block");

}
\end{lstlisting}
\par
\medskip

\par
\medskip
\begin{lstlisting}[style=out]
using
drop conn
after block
\end{lstlisting}
\par
\medskip

注意 \texttt{drop conn} 出现的位置：在里层花括号结束的那一刻，而不是在程序末尾。这就是「确定性释放」——资源什么时候还，看代码结构就能读出来。

这类拥有型资源由 \texttt{Drop} 自动收尾，不需要额外登记释放动作。结构就是释放点。 另有非所有权清理动作时可用本章的 \texttt{defer}。

\texttt{drop(\allowbreak{})\allowbreak{}} 里拿自己用 \texttt{self}。第 9 章会专门讲 \texttt{self}，现在当成「这个对象自己」就行。

\section{析构顺序：先自己的 \texttt{drop(\allowbreak{})\allowbreak{}}，再做字段}

字段会被自动递归析构——一层层往下，直到最里面：

\par
\medskip
\begin{lstlisting}[style=chali]
struct Inner: Drop
{
    String tag;
    void drop()
    {
        println("drop inner {}", self.tag);
    }
};

struct Outer: Drop
{
    Inner field;
    void drop()
    {
        println("drop outer");
    }
};

void main()
{
    Outer o;
    o.field.tag = "x";
    println("using");

}
\end{lstlisting}
\par
\medskip

\par
\medskip
\begin{lstlisting}[style=out]
using
drop outer
drop inner x
\end{lstlisting}
\par
\medskip

顺序是：先跑 \texttt{Outer} 自己的 \texttt{drop(\allowbreak{})\allowbreak{}}，再析构它的字段。

\begin{tip}{提示}
如果你写过 C++，注意这里和 C++ 是反的：C++ 先析构成员再跑析构体。ChaLi 先跑你自己的 \texttt{drop(\allowbreak{})\allowbreak{}}，之后字段自动清理。\par 这个顺序有个实际后果：\texttt{drop(\allowbreak{})\allowbreak{}} 里还能安全地用字段（\texttt{o.\allowbreak{}field} 那时还活着）。这在「关闭连接时要写日志」「释放前要刷缓冲」这类场景里很有用。
\end{tip}

\subsection*{一个必须记住的禁忌}

结构体的 \texttt{drop(\allowbreak{})\allowbreak{}} 里不要手写 \texttt{self.\allowbreak{}字段.\allowbreak{}drop(\allowbreak{})\allowbreak{}}。

因为字段本来就会被自动析构，你再手动来一次就是析构两次。这条规则没有例外：只要字段实现了 \texttt{Drop}，ChaLi 就会负责到底，你插手就是添乱。

\section{\texttt{Closeable}：用完立刻关，而且不会关两次}

有些资源不能等到「出块」才放：文件句柄、数据库连接、锁。它们要的是「用完立刻放，别占着」。

实现 \texttt{Closeable} 就是干这个的：

\par
\medskip
\begin{lstlisting}[style=chali]
struct Handle: Closeable
{
    int64 id;
    void close()
    {
        println("close {}", self.id);
    }
};

void main()
{
    {
        Handle h;
        h.id = 1;
        println("using");
    }
    println("after block");

}
\end{lstlisting}
\par
\medskip

\par
\medskip
\begin{lstlisting}[style=out]
using
close 1
after block
\end{lstlisting}
\par
\medskip

机制和 \texttt{Drop} 一样（离开作用域自动调 \texttt{close(\allowbreak{})\allowbreak{}}），但多了一条很实用的保证：

\begin{tip}{提示}
显式调用过 \texttt{close(\allowbreak{})\allowbreak{}} 的变量，不会再自动关闭。
\end{tip}

所以你可以放心地「提前关」：

\par
\medskip
\begin{lstlisting}[style=chali]
struct Handle: Closeable
{
    void close()
    {
        println("close");
    }
};

void main()
{
    {
        Handle h;
        h.close();
        // 提前关
        println("closed early");
    }
    // 这里不会再关一次
    println("after");

}
\end{lstlisting}
\par
\medskip

\par
\medskip
\begin{lstlisting}[style=out]
close
closed early
after
\end{lstlisting}
\par
\medskip

\texttt{close} 只出现了一次。 这条保证让「提前关闭」变成一件可以随手做的事——不用记「我关过了吗」，也不用写一堆标志位。

\subsection*{两个都实现的时候，顺序固定}

\par
\medskip
\begin{lstlisting}[style=out]
先 close()，再 drop()
\end{lstlisting}
\par
\medskip

不用记特例：关动作在前，销毁动作在后。

\section{堆对象：\texttt{new} 和 \texttt{drop}}

前面所有东西都活在栈上，跟着作用域走。要手动控制生命周期时，用堆：

\par
\medskip
\begin{lstlisting}[style=chali]
struct Point
{
    int64 x;
    int64 y;
};

void main()
{
    Point mut& p = heap Point;
    p.x = 3;
    p.y = 4;
    println("{},{}", p.x, p.y);
    p.drop();
    println("deleted");

}
\end{lstlisting}
\par
\medskip

\par
\medskip
\begin{lstlisting}[style=out]
3,4
deleted
\end{lstlisting}
\par
\medskip

\begin{itemize}

\item \texttt{Point mut\& p =\allowbreak{} heap Point;} —— 在堆上建一个 \texttt{Point}，\texttt{p} 是指向它的引用；

\item \texttt{p.\allowbreak{}drop(\allowbreak{})\allowbreak{};} —— 释放它。

\end{itemize}

\texttt{p} 的类型是引用（\texttt{Point mut\&}），因为堆上的东西没有名字，你只能拿着一个地址。

\texttt{new T(\allowbreak{}.\allowbreak{}.\allowbreak{}.\allowbreak{})\allowbreak{}} 也可以带实参调构造器——左边的类型决定用谁构造，那是第 9 章的内容。

\texttt{new} 还有几处等价的写法糖（语义完全一致，纯写法降级）：\texttt{heap T;} 与 \texttt{new;} 等价，\texttt{new T(\allowbreak{}实参)\allowbreak{};} 与 \texttt{new(\allowbreak{}实参)\allowbreak{};} 等价，\texttt{heap new T(\allowbreak{}实参)\allowbreak{};} 也与 \texttt{new(\allowbreak{}实参)\allowbreak{};} 等价。\texttt{heap} 是保留关键字，只在该分配位置有意义，不能用作变量名：

\par
\medskip
\begin{lstlisting}[style=chali]
Point mut& a = heap Point;             // 糖：heap T 等价于 new;
Point mut& b = new Point(3, 4);        // new T(args) 栈/堆由左侧声明决定
Point mut& c = heap new Point(5, 6);   // 糖：heap new T(args) 等价于 new(5, 6)
\end{lstlisting}
\par
\medskip

忘记 \texttt{drop} 会怎样？ 手动分配的对象不会因为引用变量离开作用域而自动释放。运行结束时会报告仍未释放的堆块；应显式 drop，或将分配封装进有明确 Drop/Closeable 契约的拥有类型。遗留分配会占用预算，长期运行可能耗尽内存。堆对象的最佳实践还是让它的生命周期和某层花括号对齐，别让 \texttt{new} 和 \texttt{drop} 离得太远。

\subsection*{编译器合并已知的重复终结}

\par
\medskip
\begin{lstlisting}[style=chali]
struct Box: Drop
{
    int64 id;
    void drop()
    {
        println("drop {}", self.id);
    }
};

void main()
{
    Box mut& b = heap Box;
    b.id = 7;
    b.drop();
    b.drop();
    println("twice ok");

}
\end{lstlisting}
\par
\medskip

\par
\medskip
\begin{lstlisting}[style=out]
drop 7
twice ok
\end{lstlisting}
\par
\medskip

输出中的 \texttt{drop} 只出现了一次。编译器能证明两条语句终结同一生命周期时，会合并它们，使析构和释放各执行一次；协议方法 \texttt{drop(\allowbreak{})\allowbreak{}} 使用同一机制。底层释放本身不幂等，跨调用等无法证明的重复释放会报告 CHL1205。

但释放之后不能再读：

\par
\medskip
\begin{lstlisting}[style=chali]
struct Point
{
    int64 x;
    int64 y;
};

void main()
{
    Point mut& p = heap Point;
    p.x = 1;
    p.drop();
    p.x = 2;

}
\end{lstlisting}
\par
\medskip

\par
\medskip
\begin{lstlisting}[style=out]
chali: [CHL1205] 借用来源已释放，不能继续访问：表达式
\end{lstlisting}
\par
\medskip

\texttt{CHL1205} 也是第 6 章 \texttt{move} 之后读源值时报的码——同一个意思：「这东西已经不在了，别再碰」。

\section{内存模型的边界}

这一节讲 ChaLi 的「安全」到底覆盖到哪儿。

会被检查的：

\begin{itemize}

\item 数组越界、整数运算溢出，都会被抓到并给出定位，不是未定义行为；

\item 悬垂引用、释放后访问、移动后访问，编译期或运行期报 \texttt{CHL1204} / \texttt{CHL1205}。

\end{itemize}

资源在哪儿：

\begin{itemize}

\item PC 上默认内存池 512 MiB（可用 \texttt{--memory-size 字节} 调整，最小 64）；

\item 嵌入式目标按板子配置，见第 22 章。

\end{itemize}

长时间运行的服务可以用 \texttt{std:\allowbreak{}:\allowbreak{}watchdog} 开「软重启」：内存申请失败、显式请求重启、越过内存水位、或开启超时后整整一个周期没 \texttt{kick(\allowbreak{})\allowbreak{}} 时，重启自身逻辑而不杀进程。

最后说清它不做什么。 ChaLi 不试图消除所有错误。它针对性地消除三类：

\begin{enumerate}

\item 空指针解引用 —— 没有 \texttt{null}，引用必须绑定真实对象（第 7 章）；

\item 悬垂引用 —— 来源作用域检查（第 7 章）；

\item 忘记释放 —— 作用域 + \texttt{Drop} / \texttt{Closeable} 自动收尾（本章）。

\end{enumerate}

剩下的逻辑错误仍然是你自己的责任。 比如「算错了平均值」「把用户的年龄存成了负数」「该关的连接关了两次但业务上不允许」——这些语言帮不了你，只能靠第 16 章的错误处理和第 18 章的测试。

\subsection*{\texttt{defer} 和 \texttt{scope(\allowbreak{}exit)\allowbreak{}}}

\texttt{defer stmt;} 登记当前作用域离开时执行的动作，\texttt{scope(\allowbreak{}exit)\allowbreak{} \{ .\allowbreak{}.\allowbreak{}.\allowbreak{} \}} 是块形式。普通结束和 \texttt{return}、\texttt{panic}、\texttt{break}、\texttt{continue} 都执行已登记动作一次；同一作用域按登记的逆序执行。

\par
\medskip
\begin{lstlisting}[style=chali]
void main()
{
    {
        defer println("first");
        scope (exit)
        {
            println("second");
        }
        println("body");
    }

}
\end{lstlisting}
\par
\medskip

\par
\medskip
\begin{lstlisting}[style=out]
body
second
first
\end{lstlisting}
\par
\medskip

\texttt{defer} 用于没有所有权关系的收尾，例如恢复状态或结束宿主操作。拥有型资源继续按 \texttt{Drop}、\texttt{Closeable} 的原规则清理。不要对同一资源再登记一份释放动作。

\section{小结}

\begin{center}
\begin{tabular}{ll}
\toprule
你会了 & 关键点 \\
\midrule
作用域 & 每对 \texttt{\{ \}} 一层；函数体和 \texttt{main} 也是作用域 \\
\texttt{Drop} & 实现它，\texttt{drop(\allowbreak{})\allowbreak{}} 在离开作用域时自动调用 \\
析构顺序 & 先自己的 \texttt{drop(\allowbreak{})\allowbreak{}}，再自动析构字段（和 C++ 相反） \\
禁忌 & 不要在 \texttt{drop(\allowbreak{})\allowbreak{}} 里手动析构字段，会双重析构 \\
\texttt{Closeable} & 离开作用域或提前 \texttt{close(\allowbreak{})\allowbreak{}}；关过就不会再关 \\
两者都有 & 顺序固定：先 \texttt{close(\allowbreak{})\allowbreak{}} 再 \texttt{drop(\allowbreak{})\allowbreak{}} \\
堆对象 & \texttt{T mut\& p =\allowbreak{} heap T;} 建，\texttt{p.\allowbreak{}drop(\allowbreak{})\allowbreak{};} 释放 \\
非幂等终结 & 编译器合并已知重复终结；其他重复释放报 CHL1205 \\
释放后访问 & \texttt{CHL1205} \\
被消除的三类错误 & 空指针、悬垂引用、忘记释放 \\
不保证的 & 逻辑错误——那是你的活儿 \\
\bottomrule
\end{tabular}
\end{center}

\section{下一步}

到此为止，「ChaLi 的骨头」这一部分讲完了：值语义、引用、资源与作用域。你现在理解了这门语言所有的取舍——为什么复制、为什么不给 \texttt{null}、为什么不用 GC。

接下来是组织代码的工具。第一个是结构体：把一堆字段和一堆方法收进同一个名字下面。

翻到\hyperref[ch:09]{第 9 章 结构体}。

\section{边界与验证}

作用域结束、显式关闭和拥有关系转移会影响资源可用性。析构协议只处理由当前对象负责的资源；裸内存与宿主句柄仍需遵守各自接口约定。

\section{习题}

前两题用于理解和基本实现，第三题检查边界，最后一题为综合挑战。先独立完成，再核对参考解答。

\begin{enumerate}

\item 资源拥有者与资源引用各负责什么？

\item 把一个文件操作的成功与失败退出路径列出来。

\item 说明 move 后为何不能继续依赖旧变量拥有该资源。

\item 挑战：为一个带 Drop 或 Closeable 协议的类型设计正常退出和提前返回的观察实验。

\end{enumerate}

\subsection*{参考解答与验收要点}

\begin{enumerate}

\item 拥有者负责按协议结束资源生命周期；引用按权限访问，并不能替代拥有者保证存活。

\item 对每条已获得资源的路径检查清理；获取失败时不能释放未成功取得的资源。

\item 拥有关系已转移，后续操作应通过接收方；原变量的可用性遵循移动检查。

\item 记录取得与释放次数，预期每份实际拥有资源按协议清理一次，不用虚构的析构顺序代替运行证据。

\end{enumerate}

\section{番外：出了作用域的穿墙术}

改编自《聊斋志异·劳山道士》，以下为编程主题的独立改写。

王生在山中学了穿墙术，师父先给他一把钥匙，说练功房开着时可以借用里面的器具。王生练了几日，觉得钥匙比器具还轻巧，回家便把它挂在门口，向妻子夸口。

妻子问那间房是否还在。王生说钥匙既是真的，门自然也该听话。他闭眼向墙走去，撞得头上一个大包。

师父闻讯，只问离山时有没有交回器具。王生说没有带走，只记住了它们的位置。师父便把借用与拥有各写一行：房间关了，记得再清楚也不能从旧位置取东西；真要带回家，必须按规矩取得自己的那一份。

王生还想求一个永不失效的法诀。师父让他先学会在离开房间时收拾自己负责的物件。妻子听了，把钥匙收进抽屉，墙仍旧好好地立在那里。

\subsection*{技术对照}

作用域结束、显式关闭和拥有关系转移会影响资源可用性。析构协议只处理由当前对象负责的资源；裸内存与宿主句柄仍需遵守各自接口约定。
